\begin{table}[htbp]
  \centering
  \caption{Jerarquía de modelos de electrodos para Tomografía de Impedancia Eléctrica.}
  \label{tab:electrode_models}
  \begin{tabular}{p{3.0cm} p{4.2cm} p{2.8cm} p{3.2cm} p{3.2cm}}
    \hline\hline
    \textbf{Modelo de Electrodo} & \textbf{Condición en Frontera ($\partial\Omega$)} & \textbf{Parámetros Requeridos} & \textbf{Ventajas Numéricas} & \textbf{Limitaciones / Singularidades} \\ \hline
    Point Electrode Model (PEM) & $\sigma \frac{\partial u}{\partial n} = \sum_{\ell} I_\ell \delta(x - x_\ell)$ & Ubicación nodal $x_\ell$ & Implementación trivial en mallas 1D/2D. & Singularidad logarítmica de potencial en 2D ($u \notin H^1$). Violación física. \\ \hline
    Shunt Model (Este trabajo) & $\sigma \frac{\partial u}{\partial n} = -\frac{I_j}{|\mathcal{E}_j|}, \quad x \in \mathcal{E}_j$ & Área de contacto $|\mathcal{E}_j|$, ancho $w=0.35$ & Elimina singularidades, $u \in H^1(\Omega)$, integración variacional exacta. & Desprecia la impedancia resistiva de contacto de la interfaz electrodo-tejido. \\ \hline
    Complete Electrode Model (CEM) & $u + z_\ell \sigma \frac{\partial u}{\partial n} = U_\ell, \quad \int_{\mathcal{E}_\ell} \sigma \frac{\partial u}{\partial n} ds = I_\ell$ & Impedancia de contacto $z_\ell$, área $|\mathcal{E}_\ell|$ & Modelo con mayor fidelidad experimental física en laboratorio y clínica. & Introduce incógnitas adicionales $U_\ell$ y acoplamiento algebraico-diferencial. \\ \hline
    \hline\hline
  \end{tabular}
\end{table}
